\chapter{Learning by Self-Play}\label{ch:learn}

Both networks are trained with Proximal Policy Optimization (PPO), an on-policy
actor--critic method, in a self-play loop: the policy being trained controls most of the
snakes in every game. This chapter states the objective, the self-play scheme, the
network, and the compute.

\section{Actor and critic}

A single network with parameters $\theta$ has two heads. Given observation $o$, it outputs
action logits, from which a masked softmax gives the policy $\pi_\theta(a\mid o)$, and a
scalar value estimate $V_\theta(o)$ of the return. The value head is what the web page
displays as the network's ``confidence'' (Chapter~\ref{ch:browser}).

\section{Advantages}

\begin{defbox}[Generalized advantage estimation]
For one seat's trajectory with rewards $r_t$ and episode boundaries $d_t$, define the TD
residual and the advantage estimate
\[
  \delta_t = r_t + \gamma\,(1-d_t)\,V(o_{t+1}) - V(o_t),
  \qquad
  \hat A_t = \sum_{l\ge 0} (\gamma\lambda)^l \prod_{m=0}^{l-1}(1-d_{t+m})\;\delta_{t+l},
\]
with $\gamma=0.99$ and $\lambda=0.95$, computed by one backward pass over the rollout. The
value targets are $\hat R_t = \hat A_t + V(o_t)$. The last step of a rollout bootstraps
from the value of the following observation.
\end{defbox}

The factor $(1-d_t)$ is where the dones of Chapter~\ref{ch:obs} matter: since every slot of
a finished game is marked done, no advantage ever leaks from one game into the next game
played in the same environment slot.

\section{The PPO objective}

Let $\pi_{\text{old}}$ be the policy that collected the rollout and
\[
  \rho_t(\theta) = \frac{\pi_\theta(a_t\mid o_t)}{\pi_{\text{old}}(a_t\mid o_t)}.
\]

\begin{defbox}[Clipped surrogate and total loss]
\[
  L^{\text{CLIP}}(\theta) = \E_t\!\left[\min\!\Bigl(\rho_t(\theta)\,\hat A_t,\;
     \clip\bigl(\rho_t(\theta),\,1-\epsilon,\,1+\epsilon\bigr)\,\hat A_t\Bigr)\right],
  \qquad \epsilon = 0.2,
\]
\[
  \mathcal{L}(\theta) = -L^{\text{CLIP}}(\theta)
    + c_v\,\E_t\!\left[\tfrac12\bigl(V_\theta(o_t)-\hat R_t\bigr)^2\right]
    - c_e\,\E_t\!\left[\Hent\bigl(\pi_\theta(\cdot\mid o_t)\bigr)\right],
\]
with $c_v=0.5$ and $c_e=0.01$. Advantages are normalised to zero mean and unit variance
over each batch.
\end{defbox}

\begin{keybox}[Why PPO here]
The game is on-policy friendly: experience is cheap (Chapter~\ref{ch:sim}), and the
data distribution shifts every iteration anyway because the opponent \emph{is} the
policy. PPO's clipping keeps each update close to the policy that gathered the data. That
matters doubly in self-play, where a large jump changes both sides of every game at
once.
\end{keybox}

\section{Self-play}

\begin{algbox}[One training iteration]
\begin{enumerate}
\item With probability $0.3$ per environment, the non-learner seats are played by a
  \emph{frozen snapshot}: one set of past weights, drawn uniformly from a pool for this
  iteration. Otherwise the learner plays every seat.
\item Run $T=64$ turns in all environments, sampling actions from the masked policy.
  Store observations (float16), masks, actions, log-probabilities, values and rewards.
\item Mark a transition \emph{valid} if its seat was played by the learner and was active
  (Chapter~\ref{ch:obs}). Compute GAE over the full rollout, then keep valid transitions
  only.
\item Optimise $\mathcal{L}$ for 4 epochs over the valid transitions in minibatches of
  4{,}096, using Adam (learning rate $2.5\times10^{-4}$, $\epsilon=10^{-5}$) and gradient
  norm clipping at $0.5$.
\item Every 10 iterations, add the current weights to the snapshot pool (the most recent
  20 are kept). Every 25--30 iterations, evaluate against the scripted bots and save a
  checkpoint.
\end{enumerate}
\end{algbox}

\begin{keybox}[Two sources of opponents]
Playing the current self is the strongest signal, and every seat yields training data,
so a two-snake game produces two transitions per turn. Playing frozen past selves guards
against \emph{forgetting}: a policy that only ever meets its current self can drift into a
strategy that loses to something it used to beat. Seats a snapshot plays generate no
training data.
\end{keybox}

Each environment re-draws its opponent type whenever its game ends. In the duel run,
$512$ environments $\times$ $64$ turns gives $32{,}768$ game turns per iteration and
about $55{,}000$ valid learner transitions. The arena run uses $256$ environments of six
snakes each.

\section{The network}

\begin{center}
\begin{tikzpicture}[node distance=0.42cm and 0.38cm,
    layer/.style={node box, minimum height=1.25cm, text width=1.55cm},
    hd/.style={node fill, text width=1.25cm}]
  \node[node soft, text width=1.65cm, minimum height=1.25cm] (in) {input\\[1pt]\scriptsize $C{+}1$ planes\\$H\times W$};
  \node[layer, right=of in] (c1) {conv $3{\times}3$\\\scriptsize 32, ReLU};
  \node[layer, right=of c1] (c2) {conv $3{\times}3$\\\scriptsize 64, ReLU};
  \node[layer, right=of c2] (c3) {conv $3{\times}3$\\\scriptsize 64, ReLU};
  \node[layer, right=of c3] (fc) {linear\\\scriptsize $64HW\!\to\!h$, ReLU};
  \node[hd, above right=0.05cm and 0.5cm of fc] (pi) {policy\\\scriptsize 4 logits};
  \node[hd, below right=0.05cm and 0.5cm of fc] (v) {value\\\scriptsize 1 scalar};
  \draw[flow] (in) -- (c1);
  \draw[flow] (c1) -- (c2);
  \draw[flow] (c2) -- (c3);
  \draw[flow] (c3) -- (fc);
  \draw[flow] (fc.east) -- (pi.west);
  \draw[flow] (fc.east) -- (v.west);
\end{tikzpicture}
\captionof{figure}{The actor--critic. All convolutions use padding 1, so spatial size is
preserved; the extra input plane is the constant ones plane that reveals walls through
zero padding.}
\label{fig:net}
\end{center}

\begin{center}
\begin{tabular}{@{}lccc@{}}
\toprule
& \textbf{Input} $C{+}1,\,H{\times}W$ & \textbf{Hidden} $h$ & \textbf{Parameters}\\
\midrule
Duel network & $9,\ 11\times11$ & 256 & 2{,}042{,}053\\
Arena network & $10,\ 17\times17$ (window) & 128 & 2{,}426{,}597\\
\bottomrule
\end{tabular}
\end{center}

Almost all parameters sit in the first linear layer, $64HW\cdot h$, which is why the arena
network halves $h$ to compensate for its larger window. A fully convolutional design
would remove the dependence on $H\times W$ altogether. The window observation achieves
the same independence from board size while keeping this simple architecture.

\section{Compute}

Both networks trained for 2.5 hours on an Apple M3 Max, using PyTorch's Metal (MPS)
backend for the network and the CPU for the simulator.

\begin{center}
\begin{tabular}{@{}lcc@{}}
\toprule
& \textbf{Seconds per iteration} & \textbf{Relative}\\
\midrule
CPU, 12 threads & $\approx 70$ & $1\times$\\
Apple GPU (MPS) & $\approx 5$--$6$ & $\approx 10\times$\\
\bottomrule
\end{tabular}
\end{center}

\begin{warnbox}[The GPU mattered more than expected]
The network is small, so CPU training looked plausible. In practice the PPO update,
four epochs over $\sim55{,}000$ samples through a $2$M-parameter network, dominated each
iteration, and moving it to the GPU was the difference between a 2.5-hour run and a
day-long one. The simulator stayed on the CPU throughout.
\end{warnbox}
